\section{Combinatorics and the degenerate affine Hecke algebra}\label{sec:deg_hecke}

\subsection{Root ideals and reverse Hessenberg functions}\label{ss:root_ideals_hess_funcs}

Let $n\geq2$ and set $[n]=\{1,\dots,n\}$. Let $G=\mathrm{GL}_n(\C)$, with Borel subgroup $B$
of upper triangular matrices and maximal torus $T$ of diagonal matrices.
Write $\mathfrak g=\mathfrak{gl}_n(\C)$ and $\mathfrak t=\operatorname{Lie}(T)$
for the Lie algebras of $G$ and $T$, respectively.
Let $\epsilon_i\in\mathfrak t^*$ be the $i$-th coordinate character of $T$.
The root system of $(G,T)$ and its positive roots are
\[
  \Delta=\{\epsilon_i-\epsilon_j\mid i\neq j\},
  \qquad
  \Delta^+=\{\epsilon_i-\epsilon_j\mid i<j\}.
\]
A subset $\Psi\subset\Delta^+$ is a \emph{root ideal} if
$(\Psi+\Delta^+)\cap\Delta^+\subset\Psi$.
Any root ideal $\Psi$ is uniquely encoded by a
\emph{reverse Hessenberg function}, that is, a weakly increasing map
$\rfunc:[n]\to\{0,\dots,n-1\}$ satisfying $\rfunc(j)<j$ for all $j\in[n]$.
The bijection between root ideals and reverse Hessenberg functions is given by
\[
  \Psi_{\rfunc}
  =\{\epsilon_i-\epsilon_j\in\Delta^+\mid i\leq\rfunc(j)\},
  \qquad
  \rfunc_\Psi(j)
  =\max\bigl(\{i\mid\epsilon_i-\epsilon_j\in\Psi\}\cup\{0\}\bigr).
\]
The Weyl group $W=N_G(T)/T$ is identified with the symmetric group $\mathfrak S_n$,
where the simple reflection $s_i=(i,i+1)$ interchanges $i$ and $i+1$.

\begin{remark}[Hessenberg functions and reverse Hessenberg functions]\label{rem:hessenberg_comparison}
In combinatorics and the geometry of Hessenberg varieties, it is more common to encode a root ideal $\Psi\subset\Delta^+$ by a \emph{Hessenberg function}, namely a weakly increasing map $\hfunc:[n]\to[n]$ satisfying $\hfunc(i)\geq i$ for all $i\in[n]$. A Hessenberg function $\hfunc$ determines a root ideal by $\Psi_\hfunc=\{\epsilon_i-\epsilon_j\in\Delta^+\mid j>\hfunc(i)\}$. The bijection between Hessenberg functions $\hfunc$ and reverse Hessenberg functions $\rfunc$ corresponding to the same root ideal is given by the equivalence $i\leq\rfunc(j)\iff j>\hfunc(i)$, or explicitly by
\[
  \rfunc(j)=\max\bigl(\{i\in[n]\mid \hfunc(i)<j\}\cup\{0\}\bigr),
  \qquad
  \hfunc(i)=\max\{j\in[n]\mid\rfunc(j)<i\}.
\]
Although Hessenberg functions are more standard in the combinatorics literature, we follow Kamnitzer \cite[\S2.1]{kamnitzer2024fixed} in using reverse Hessenberg functions throughout this paper, as they allow us to write the defining lattice condition of $\SX_{\rfunc}$ and the rank of each stage of the projective bundle tower directly in terms of $\rfunc(j)$ (see Section~\ref{sec_xpsi}).
\end{remark}

\subsection{The degenerate affine Hecke algebra and polynomial operators}\label{ss:deg_hecke_algs}

Let $\epsilon$ be a formal parameter of degree two.
We recall the degenerate affine Hecke algebra (cf.\ Ginzburg \cite[\S 12]{ginzburg1998geometric}).

\begin{definition}\label{def:deg-hecke}
The \emph{degenerate affine Hecke algebra} $\Hdeg$ is the associative algebra over $\Z[\epsilon]$
generated by pairwise commuting elements $x_1,\dots,x_n$ and simple reflections $s_1,\dots,s_{n-1}$
satisfying the Coxeter relations of $\mathfrak S_n$ together with the cross relations
\[
  s_jx_i-x_{s_j(i)}s_j=
  \begin{cases}
    -\epsilon & \text{if } i=j,\\
    \epsilon & \text{if } i=j+1,\\
    0 & \text{if } i\neq j,j+1,
  \end{cases}
  \qquad (1\leq i\leq n,\ 1\leq j<n).
\]
The algebra $\Hdeg$ is graded by $\deg s_j=0$ and $\deg x_i=\deg\epsilon=2$.
\end{definition}

Let $\Lambda$ be a commutative $\Z[\epsilon]$-algebra and set $R=\Lambda[x_1,\dots,x_n]$.
Let $\tau_i$ denote the $\Lambda$-algebra automorphism of $R$ interchanging
$x_i$ and $x_{i+1}$ and fixing all other variables.
We define the divided difference operator $\partial_i$ and the
Demazure--Lusztig type operator $\BS_i$ on $R$ by
\[
  \partial_i f=\frac{f-\tau_i f}{x_i-x_{i+1}},
  \qquad
  \BS_i f=\tau_i f-\epsilon\partial_i f.
\]
The divided difference operators satisfy the nil-Hecke relations and the
twisted Leibniz rule
\begin{equation}\label{eq:divided_difference_leibniz}
  \partial_i(fg)=(\partial_i f)g+(\tau_i f)\partial_i g.
\end{equation}
The operators $\BS_i$ satisfy the Coxeter relations of $\mathfrak S_n$, and together
with multiplication by $x_i$, they define the polynomial representation
\begin{equation}\label{eq:deg-hecke-action}
  s_i\longmapsto\BS_i,\qquad
  x_i\longmapsto x_i
\end{equation}
of $\Hdeg$ on $R$ (see \cite[Proposition~12.2]{ginzburg1998geometric}).
Note that $\BS_i$ can be rewritten in the form
\begin{equation}\label{eq:BS-expansion}
 \BS_i=\id+(x_{i+1}-x_i-\epsilon)\partial_i.
\end{equation}
When $R$ is graded with $\deg\epsilon=2$ and $\deg x_i=2$, the
operator $\partial_i$ has degree $-2$ and $\BS_i$ has degree zero.

\begin{lemma}\label{lem:polynomial_faithful}
The polynomial representation of $\Hdeg$ on
$\Z[x_1,\dots,x_n,\epsilon]$ is faithful.
\end{lemma}

\begin{proof}
For $w\in W$, let $s_w$ be the product of simple generators along a reduced expression.
By the cross relations, every element of $\Hdeg$ can be written as
$\sum_{w\in W}p_w s_w$ with $p_w\in\Z[x_1,\dots,x_n,\epsilon]$.
If this element acts by zero on $\Z[x_1,\dots,x_n,\epsilon]$, then its image in
$\C\otimes_{\Z}\Hdeg$ acts by zero on $\C[x_1,\dots,x_n,\epsilon]$ and is therefore zero
by \cite[Proposition~12.2(ii)]{ginzburg1998geometric}.
The complex PBW decomposition \cite[Definition~12.1(i)]{ginzburg1998geometric}
then gives $p_w=0$ for every $w\in W$, so the element is zero in $\Hdeg$.
\end{proof}

\subsection{Polynomials for the cohomology presentation}\label{ss:lemmas_polynomials}

We now establish polynomial identities that express the relations arising from the
projective bundle formula for the tower defining $\SX_{\rfunc}$ (see Section~\ref{sec_xpsi}).

Let $f(z)\in\Lambda[z]$ be a monic polynomial of degree $n$.
For $0\leq j<k\leq n$, define the polynomial $P_{k,j}(f)\in R$ by
\[
  P_{k,j}(f)
  =\partial_{k-1}\cdots\partial_{j+1}
       \BS_j\cdots\BS_1 f(x_1),
\]
where empty products of operators are understood to be the identity.
For a reverse Hessenberg function $\rfunc$, set
\[
  D_k(f)=P_{k,\rfunc(k)}(f),\qquad
  I_{\rfunc}^{[k]}(f)=(D_1(f),\dots,D_k(f))\subset R,
  \qquad I_{\rfunc}^{[0]}(f)=0,
\]
and write $I_{\rfunc}(f)=I_{\rfunc}^{[n]}(f)$.

\begin{lemma}\label{lem:polynomial_recurrence}
The polynomial $P_{k,j}(f)$ belongs to $\Lambda[x_1,\dots,x_k]$ and is monic
of degree $n-k+1+j$ in $x_k$.
Moreover, for all $1\leq j<k$, it satisfies the recurrence
\begin{equation}\label{eq:triangular_polynomial_recurrence}
  P_{k,j}(f)
  =(x_k-x_j-\epsilon)P_{k,j-1}(f)+P_{k-1,j-1}(f).
\end{equation}
\end{lemma}

\begin{proof}
If $g\in\Lambda[x_1,\dots,x_i]$ is monic of degree $d>0$ in $x_i$,
then $\BS_i g$ and $\partial_i g$ are monic in $x_{i+1}$ of degrees $d$ and $d-1$,
respectively.
In the definition of $P_{k,j}(f)$, we apply $j$ Demazure--Lusztig type operators to $f(x_1)$,
followed by $k-j-1$ divided differences.
Thus $P_{k,j}(f)\in\Lambda[x_1,\dots,x_k]$ is monic of degree $n-k+1+j$ in $x_k$.

For fixed $j$, the case $k=j+1$ of
\eqref{eq:triangular_polynomial_recurrence} follows from
\eqref{eq:BS-expansion} applied to $P_{j,j-1}(f)$.
For $k>j+1$, we apply $\partial_{k-1}$ to the recurrence for $k-1$.
Since $P_{k-2,j-1}(f)\in\Lambda[x_1,\dots,x_{k-2}]$ is independent of $x_{k-1}$ and $x_k$,
it is annihilated by $\partial_{k-1}$, and the Leibniz rule
\eqref{eq:divided_difference_leibniz} gives the result, using
$\partial_{k-1}(x_{k-1}-x_j-\epsilon)=1$ and
$\tau_{k-1}(x_{k-1}-x_j-\epsilon)=x_k-x_j-\epsilon$.
\end{proof}

We introduce the reversed polynomials in an auxiliary variable $u$:
\[
  F(u)=u^nf(u^{-1}),\qquad
  d_{k,j}=n-k+1+j,\qquad
  \widehat P_{k,j}(u)
    =u^{d_{k,j}}P_{k,j}(f)\big|_{x_k=u^{-1}}.
\]
By Lemma~\ref{lem:polynomial_recurrence}, $\widehat P_{k,j}(u)\in \Lambda[x_1,\dots,x_{k-1}][u]$
is a polynomial in $u$ of degree at most $d_{k,j}$ with constant term $1$.
The following lemma computes $\widehat P_{k,j}(u)$ modulo $I_{\rfunc}^{[k-1]}(f)$.

\begin{lemma}\label{lem:operator_chern_relation}
Let $1\leq k\leq n$ and $\rfunc(k)\leq j<k$.
\begin{enumerate}
  \item\label{item:Pk}
    $P_{k,j}(f)\in I_{\rfunc}^{[k]}(f)$.
  \item\label{item:Qk}
    In $(R/I_{\rfunc}^{[k-1]}(f))\llbracket u\rrbracket$,
    \begin{equation}\label{eq:operator_chern_series}
      \widehat P_{k,j}(u)
      =
      F(u)\,
      \frac{\prod_{i=1}^{j}(1-(x_i+\epsilon)u)}
           {\prod_{i=1}^{k-1}(1-x_i u)}.
    \end{equation}
\end{enumerate}
\end{lemma}

\begin{proof}
We prove (i) by induction on $k$. The base case $k=1$ is
$P_{1,0}(f)=D_1(f)\in I_{\rfunc}^{[1]}(f)$.
For $k\geq 2$, the boundary case $j=\rfunc(k)$ holds by definition:
$P_{k,\rfunc(k)}(f)=D_k(f)\in I_{\rfunc}^{[k]}(f)$.
For $j>\rfunc(k)$, the monotonicity of $\rfunc$ implies $j-1\geq\rfunc(k)\geq\rfunc(k-1)$,
so $P_{k-1,j-1}(f)\in I_{\rfunc}^{[k-1]}(f)$ by induction on $k$.
The recurrence \eqref{eq:triangular_polynomial_recurrence} then gives
$P_{k,j}(f)\in I_{\rfunc}^{[k]}(f)$ by induction on $j$.

We prove (ii) by induction on $k$. The base case $k=1$ is $\widehat P_{1,0}(u)=F(u)$.
For $k\geq 2$, we distinguish two cases.

If $j\leq k-2$, then $P_{k,j}(f)=\partial_{k-1}P_{k-1,j}(f)$.
Since $P_{k-1,j}(f)\in\Lambda[x_1,\dots,x_{k-1}]$ is independent of $x_k$, applying $\tau_{k-1}$ substitutes $x_k$ for $x_{k-1}$, so
\[
  u^{d_{k-1,j}}(\tau_{k-1}P_{k-1,j}(f))\big|_{x_k=u^{-1}}=\widehat P_{k-1,j}(u).
\]
Set $d=d_{k,j}$, so that $d_{k-1,j}=d+1$.
Expanding the divided difference and substituting $x_k=u^{-1}$ gives
\[
  (1-x_{k-1}u)\widehat P_{k,j}(u)
  =\widehat P_{k-1,j}(u)-u^{d+1}P_{k-1,j}(f).
\]
Since $j\geq\rfunc(k)\geq\rfunc(k-1)$, the second term vanishes modulo
$I_{\rfunc}^{[k-1]}(f)$ by part~(i).
Applying the induction hypothesis to $\widehat P_{k-1,j}(u)$ and dividing by
$1-x_{k-1}u$ yields \eqref{eq:operator_chern_series}.

If $j=k-1$, then $P_{k,k-1}(f)=\BS_{k-1}P_{k-1,k-2}(f)$.
Since $d_{k,k-1}=d_{k-1,k-2}=n$, applying the same substitution to
$P_{k-1,k-2}(f)\in\Lambda[x_1,\dots,x_{k-1}]$ gives
$u^n(\tau_{k-1}P_{k-1,k-2}(f))\big|_{x_k=u^{-1}}=\widehat P_{k-1,k-2}(u)$.
Using \eqref{eq:BS-expansion} and setting $x_k=u^{-1}$, we find
\[
  (1-x_{k-1}u)\widehat P_{k,k-1}(u)
  =(1-(x_{k-1}+\epsilon)u)\widehat P_{k-1,k-2}(u)
  +\epsilon u^{n+1}P_{k-1,k-2}(f).
\]
Since $k-2\geq\rfunc(k-1)$ by the condition $\rfunc(k-1)<k-1$, the second term
again vanishes modulo $I_{\rfunc}^{[k-1]}(f)$ by part~(i).
Applying the induction hypothesis to $\widehat P_{k-1,k-2}(u)$ and dividing by
$1-x_{k-1}u$ gives \eqref{eq:operator_chern_series}.
\end{proof}

In Subsection~\ref{ss:equiv_coho_xpsi}, we identify $\widehat P_{k,\rfunc(k)}(u)$,
modulo the preceding relations, with the Chern polynomial at the $k$-th stage of the projective bundle tower.

\subsection{Simple-reflection symmetries of root ideals}

The geometric realization of the operators $\BS_i$ in Section~\ref{sec:geometric_hecke}
requires the reflection $s_m$ to preserve the set $\Psi\cup(-\Psi)$.
In terms of the reverse Hessenberg function, this symmetry is characterized as follows.

\begin{lemma}\label{lem:symmetry_of_root_ideal}
Let $\Psi=\Psi_{\rfunc}$ and $1\leq m<n$.
The set $\Psi\cup(-\Psi)$ is $s_m$-invariant if and only if
$\rfunc(k)\neq m$ for all $k\geq m+2$ and one of the following two
mutually exclusive conditions holds:
\begin{enumerate}
  \item $\rfunc(m)=m-1$ and $\rfunc(m+1)=m$;\label{item:case1A}
  \item $\rfunc(m)=\rfunc(m+1)$.\label{item:case1B}
\end{enumerate}
In case~\ref{item:case1A}, $\epsilon_m-\epsilon_{m+1}\in\Psi$;
in case~\ref{item:case1B}, $\epsilon_m-\epsilon_{m+1}\notin\Psi$.
\end{lemma}

\begin{proof}
Aside from the simple root $\epsilon_m-\epsilon_{m+1}$, whose sign is reversed,
$s_m$ interchanges the roots in each of the pairs
\[
  \{\epsilon_i-\epsilon_m,\epsilon_i-\epsilon_{m+1}\}\quad(i<m),
  \qquad
  \{\epsilon_m-\epsilon_k,\epsilon_{m+1}-\epsilon_k\}\quad(k>m+1).
\]
For the first family, membership in $\Psi$ agrees within each pair
if and only if $\min(\rfunc(m+1),m-1)=\rfunc(m)$.
Since $\rfunc(m)\leq\rfunc(m+1)\leq m$, this equality is equivalent to
either (i) $\rfunc(m+1)=m$ and $\rfunc(m)=m-1$, or (ii) $\rfunc(m+1)=\rfunc(m)\le m-1$.
For the second family, membership agrees within each pair if and only if $\rfunc(k)\neq m$.
All remaining positive roots are fixed by $s_m$.
Since $\Psi\cup(-\Psi)$ is symmetric under negation, this proves the claim.
\end{proof}
